\begin{helper}
If \(n\ge 1\) and the huge cutoff satisfies
\[
100\le \noderef{TwoBiteHugeCutoff}(n),
\]
then
\[
2\le (\log n)^2.
\]
\end{helper}
\begin{proof}
Write \(L=\log n\).  Since the cutoff is positive, its denominator
\(\log L\) must be positive: the numerator \(\sqrt{nL}\) is nonnegative, so a
nonpositive denominator would make
\(\sqrt{nL}/\log L=\noderef{TwoBiteHugeCutoff}(n)\) nonpositive.  Hence
\(\log L>0\), and therefore \(L>1\).  In particular \(n>e>2\), so \(n\ge3\).

We first get a fixed positive lower bound for \(\log L\).  From \(n\ge3\),
\(L=\log n\ge\log3\).  The standard numerical exponential estimates give
\(e<14/5\) and \(e^{1/19}\le19/18\), hence
\[
e^{20/19}=e\,e^{1/19}<3.
\]
Thus \(20/19\le\log3\le L\).  Also
\[
e^{1/20}\le \frac1{1-1/20}=\frac{20}{19},
\]
so \(e^{1/20}\le L\), and consequently
\[
\frac1{20}\le \log L.
\]

Multiplying \(100\le \sqrt{nL}/\log L\) by the positive denominator
\(\log L\) gives
\[
100\log L\le\sqrt{nL}.
\]
The previous lower bound on \(\log L\) yields \(5\le\sqrt{nL}\), and squaring
gives
\[
25\le nL. \tag{1}
\]

Suppose for contradiction that \(L^2<2\).  Since \(L>0\), this implies
\(L<2\).  Therefore \(n<e^2<9\), again by the standard bound \(e<3\), so
\(n\le8\).  Combining \(n\le8\) with \(L<2\) gives \(nL<16\), contradicting
(1).  Hence \(2\le L^2=(\log n)^2\), as claimed.
\end{proof}
